% -*- coding: utf-8 -*-
\chapter{绪论}
\label{chap:intro}

\section{研究背景与低空智联网业务演进}
\label{sec:background}

\subsection{低空经济与低空智联网的国家战略定位}
随着低空空域管理改革的纵深推进与通用航空技术的跨越式发展，以无人机（Unmanned Aerial Vehicle, UAV）为核心主体的低空经济已上升为国家战略性新兴产业。低空智联网（Low-Altitude Internet of Things, LA-IoT）作为支撑低空飞行器规模化运行的关键数字化基础设施，承担着空域安全监管、飞行器协同组网、超视距控制指令下达以及宽带载荷数据回传等多重任务。

从产业形态看，低空经济并非单一行业，而是由低空制造、低空飞行、低空保障与综合服务四个板块构成的复合生态。其典型应用场景包括城市末端物流配送、电力与油气长输管线巡检、农林植保作业、应急救援与医疗物资投送、城市空中交通（Urban Air Mobility, UAM）以及低空文旅观光等。这些场景对低空智联网提出的需求并不一致：物流配送要求高可靠的指令链路与精确的航路保持；管线巡检要求对沿线设施的高分辨率成像与缺陷识别；应急救援要求在有遮挡、强干扰的复杂地形下维持通信不中断；城市空中交通则要求对高密度飞行器流进行实时冲突探测与解脱。

需求的多样性直接决定了低空智联网不能沿用单一功能的传统架构。若按功能分别建设通信网、监视网与导航网，则在同一片空域内将出现多套射频前端、多组天线阵列与多份频谱占用，站址资源与能耗开销成倍增长，而频谱在低频段已高度紧张。这一矛盾在低空场景中尤为突出：低空飞行器数量大、密度高、更新快，若每一种新业务都需要新建一套专用基础设施，则基础设施建设速度将永远滞后于业务增长。

\subsection{低空监视的现实困境}
在低空空域（飞行高度通常在 1000 米以下）环境下，飞行环境呈现出三维立体、障碍密集与动态异构的特征。低空飞行器种类繁多，涵盖消费级航拍无人机、行业巡检多旋翼平台以及垂直起降电动飞行器（eVTOL）。这些目标普遍具有低空飞行、慢速巡航及雷达散射截面积较小的“低慢小”（Low-altitude, Slow-speed, Small-size, LSS）物理属性。

“低慢小”属性给传统监视手段带来了三重困难。其一是探测距离受限：雷达散射截面积小意味着回波功率微弱，在四次方衰减规律下，可探测距离被显著压缩。其二是杂波抑制困难：低空目标的多普勒频率与地物、植被、建筑物的杂波谱高度重叠，传统动目标显示（Moving Target Indication, MTI）滤波器难以在不损失目标能量的前提下剔除杂波。其三是遮挡与多径：城市低空环境中建筑物密集，地面雷达视线被频繁遮挡，形成大量探测盲区，而多径反射又会产生虚假点迹。

传统的地面二次雷达监视网络受限于地面建筑物理遮挡，存在严重的探测盲区与多径地杂波干扰；而若单独建设高密度专用雷达基站，将面临高昂的站址资源成本与频谱匮乏约束。这两条路径各自的困境，构成了本文研究的现实起点。

\subsection{通感一体化带来的技术契机}
在此背景下，通感一体化（Integrated Sensing and Communication, ISAC）技术应运而生。ISAC 旨在打破通信与雷达探测各自独立的传统烟囱式系统架构，通过在发射端共享射频前端硬件、天线阵列与无线频谱资源，利用联合设计的基带信号同时实现对空域目标的高精度环境感知与高速可靠无线通信传输。

将 ISAC 技术深度融合至低空无人机系统，可带来三个层面的收益。在基础设施层面，ISAC 能够复用现有移动通信蜂窝网络的广域覆盖优势，无需为感知功能单独选址建站，大幅降低部署成本与能耗开销。在频谱层面，ISAC 在同一段频谱上同时承载通信与感知，频谱效率相对于功能分离方案近似翻倍，这对低频段频谱已高度紧张的现实具有直接价值。在系统能力层面，感知获得的空域目标位置与运动状态可反哺通信的动态波束对准与链路资源调度，而通信网络的高密度基站又可提供多点协同的感知视角，二者形成正向增强而非简单叠加。

正是这一“一加一大于二”的潜力，使得 ISAC 被视为构建具备全域感知与自适应通信能力的低空泛在立体智联网络的核心使能技术。然而，潜力的兑现并非无条件：通信与感知在波形、功率、波束与子载波上的资源竞争是客观存在的物理事实，如何在资源受限前提下定量刻画并优化这一权衡，是本文要回答的核心问题。

\section{无人机通感一体化核心挑战与理论矛盾}
\label{sec:challenge}

将通感一体化技术引入高动态低空无人机场景，面临诸多基础理论挑战与工程实现矛盾：

\begin{enumerate}
  \item \textbf{波形多维时频空域资源博弈矛盾}：雷达探测需要恒模、低峰均功率比（PAPR）且具备理想二维模糊函数的探测波形以保障远距离探测灵敏度；而多载波宽带数字通信追求高阶正交幅度调制（QAM）以最大化频谱效率与香农信道容量。两者在多载波子载波功率分配与相位调制上存在天然的物理竞争关系。低 PAPR 的恒模波形有利于功率放大器工作在饱和区从而提高能效，但恒模约束会压缩星座的自由度，降低可达速率；反之，高阶 QAM 带来的幅度起伏会显著抬高峰均比，迫使功放回退，实际发射功率下降，反过来削弱探测距离。

  \item \textbf{高动态低空多普勒频移瓶颈}：无人机在巡航与机动飞行过程中，空地视距（LoS）链路伴随显著的时变多普勒频偏。在毫米波等高频通信体制下，多普勒扩展严重破坏传统正交频分复用（OFDM）的正交性，产生剧烈的载波间干扰（ICI），恶化通信误码率并引发雷达虚假目标模糊。更棘手的是，多普勒频偏同时携带目标径向速度信息，是感知的重要观测量；抑制多普勒以保通信与利用多普勒以做感知，二者在信号处理层面构成直接对立。

  \item \textbf{空地非线性三维轨迹与波束赋形高维耦合}：无人机空间机动轨迹直接改变空地信道几何大尺度衰减与阵列空间导向矢量，进而影响多目标感知克拉美罗界（CRB）矩阵与多用户通信信干噪比（SINR）。这一多变量非凸耦合使得全局优化求解极具挑战。轨迹变量与波束变量在目标函数中以“距离的四次方反比”与“导向矢量的二次型”两种完全不同的方式耦合，前者关于位置强非线性，后者关于角度强非线性，二者的联合优化问题既非凸也无解析解。

  \item \textbf{感知性能度量的选择困境}：感知性能可用检测概率、估计均方误差、克拉美罗界、互信息等多种指标刻画。检测概率依赖具体的检测器与门限设置，估计均方误差依赖具体的估计算法，二者都难以作为与算法无关的系统级设计指标。克拉美罗界作为无偏估计的方差下界，与具体估计算法无关，因而更适合作为优化目标；但它关于发射协方差矩阵呈分式结构，给优化带来额外困难。如何在“指标可优化”与“指标可解释”之间取得平衡，是建模阶段必须先解决的问题。

  \item \textbf{评价基准的公平性缺失}：现有大量工作以通信速率与感知精度的加权和作为目标函数，通过调节权重生成性能折中曲线。然而加权和的两个分量量纲不同（bit/s/Hz 与 rad$^{2}$），权重本身缺乏物理含义；更重要的是，权重趋于零时加权和会退化为单一功能优化，此时另一功能的性能可能发散至非物理的奇点。这使得不同文献报告的“折中增益”之间难以横向比较，也无法回答工程上真正关心的问题：在给定通信需求下，感知精度究竟能做到多好。
\end{enumerate}

上述五项挑战中，前两项属于信号与波形层面，第三项属于优化与资源分配层面，第四、五项属于建模与评价层面。本文不试图同时解决全部问题，而是聚焦于第三、四、五项所指向的核心矛盾：\textbf{在通信需求被完全满足的前提下，如何把剩余资源最优地用于感知，并给出可横向比较的公平评价}。这一聚焦使得本文能够在一个边界清晰的框架内给出可证明、可复现、可验证的结论。

\section{国内外研究现状}
\label{sec:related}

近年来，国内外学术界与工业界围绕通感一体化波形优化、波束赋形以及空地协同调度开展了广泛研究。本节按四条线索梳理：波形与调制体制、波束赋形与资源分配、感知性能度量、轨迹优化与机动协同。

\subsection{通感一体化波形与调制体制}
在通感一体化波形设计方面，熊等人系统梳理了雷达通信一体化波形的发展脉络与调制体制\cite{xiong2022overview}，将现有方案归纳为以通信为中心的波形设计、以雷达为中心的波形设计与联合设计三条路线，并指出联合设计路线在性能上最具潜力但复杂度最高。罗等人从信息论限与基础波形结构出发，深入总结了多载波波形设计面临的权衡关系\cite{luo2024integrated}，指出多载波体制在感知分辨率上的天然优势与其高峰均比代价构成一对基本矛盾。

针对多载波系统峰均比过高的瓶颈，顾等人提出了面向无人机通信感知一体化的低 PAPR 联合波形设计方案\cite{gu2024papr}；尹等人则探索了保频谱效率的低峰均比波形构造策略\cite{yin2026papr}；吴等人进一步提出了基于符号级预编码的多用户 MIMO-OFDM 低峰均比发射方案\cite{wu2025symbol}。这三项工作从不同角度缓解了峰均比与频谱效率的对立，但都未涉及感知精度与峰均比的定量权衡关系。

在抗高多普勒先进调制体制方面，万等人针对毫米波高机动无人机场景，提出了正交线性调频多分复用（OCDM）一体化波形，利用菲涅尔变换域的能量聚集特性显著增强了高频多普勒鲁棒性\cite{wan2024orthogonal}。此外，霍等人基于检测概率最大化准则优化了 MIMO-OFDM 探测波形\cite{huo2024waveform}；张等人针对特殊目标特征探索了专用一体化微多普勒波形\cite{zhang2025waveform}；仲等人设计了恒模约束下的通感一体波形流形优化算法\cite{zhong2022constant}；赵等人则针对杂波干扰背景提出了隐蔽一体化信号传输方案\cite{zhao2025covert}。

上述波形层面工作的共同特点是：以单一功能的性能指标（速率、检测概率或峰均比）为主要优化对象，另一功能作为约束或事后评估。这一定位使得波形研究难以回答“在给定速率下感知精度上限是多少”这类系统级问题。

\subsection{波束赋形与资源分配}
在空地资源分配与机动轨迹协同方面，付等人研究了无人机辅助通感系统的联合发射波束赋形与航迹规划问题\cite{fu2024joint}；左等人探讨了非理想信道状态信息下的鲁棒波形优化方法\cite{zuo2025robust}；齐等人引入深度学习框架实现了数据驱动的波束赋形自适应优化\cite{qi2023based}；周等人分析了无人机计算、通信与雷达感知三维能耗与精度的权衡界限\cite{zhou2024trade}；张等人则系统评估了基于智能超表面辅助通感系统的互信息量界限\cite{zhang2022joint}；于等人针对高动态车联网毫米波环境分析了子载波分类调制的性能特征\cite{yu2024integrated}。

这一线索的工作在方法论上有两个共同取向。其一是采用半正定松弛（Semidefinite Relaxation, SDR）处理非凸的波束赋形问题：把关于波束矩阵的二次型约束松弛为关于协方差矩阵的线性矩阵不等式约束，求解后再通过高斯随机化恢复秩一解。该路线理论完备、性能接近最优，但存在两个实践障碍：随机化恢复可能破坏原有的速率约束可行性，需要额外的可行性修复步骤；且内点法求解半定规划的复杂度随天线数呈高次增长，在大规模阵列下计算代价可观。

其二是采用加权和作为目标函数，通过扫描权重生成性能折中曲线。该取向的问题已在第~\ref{sec:challenge} 节第四、五项挑战中说明：权重缺乏物理含义，且权重极值处会出现非物理的退化。

\subsection{感知性能度量与克拉美罗界}
克拉美罗界作为无偏估计方差的下界，是通感一体化研究中最常用的感知性能度量。对于均匀线阵的角度估计问题，经典结论给出 CRB 与阵元数的三次方成反比。这一结论在雷达文献中被广泛引用，但其成立依赖于一个容易被忽略的前提：信噪比按单阵元定义。

当回波幅度参数吸收了全部阵列增益（即把阵列增益计入等效散射系数）时，CRB 对阵元数的标度会从三次方降为平方。这一区别在系统级优化中具有实质影响：若误用三次方标度，会高估增加阵元带来的感知收益，从而在资源分配中作出偏离最优的决策。本文在第三章将对此给出严格推导，并在第四章用蒙特卡洛实验予以确认。

除 CRB 之外，互信息、检测概率与模糊函数主瓣宽度也被用作感知度量。互信息与 CRB 在高信噪比下近似等价，但在低信噪比下会给出不同排序；检测概率依赖检测器设计；模糊函数主瓣宽度只反映分辨率而不反映精度。综合比较后，本文选择 CRB 作为优化目标，理由是其与估计算法无关、可解析求导、且在高信噪比区间与估计均方误差吻合良好。

\subsection{轨迹优化与机动协同}
无人机轨迹优化是通感一体化区别于固定基站场景的独特问题。轨迹通过两条路径影响系统性能：一是改变空地距离，从而改变大尺度路径损耗；二是改变目标相对阵列的角度，从而改变导向矢量与感知几何。前者对通信与感知同向有利（距离近则两者都好），后者则可能产生冲突（有利于通信的角度未必有利于感知）。

现有工作多采用“分块坐标下降 + 凸近似”的路线：把轨迹变量分块，每块在给定其他块的前提下作凸近似（如连续凸近似、泰勒展开下界），迭代求解。该路线收敛性有保证，但收敛速度与近似紧致度依赖问题结构，且每轮都需要求解一个凸子问题，总体计算量较大。另一类工作采用深度强化学习直接学习轨迹策略，避免了对问题结构的建模，但可解释性与可复现性较弱，且训练成本高。

本文选择序列二次规划（SLSQP）求解轨迹子问题，原因是问题维度较小（两个无人机、三十个时隙共 120 维）、约束结构简单（仅速度范数约束）、且目标函数梯度可解析求出。这一选择牺牲了理论上的全局最优性，换取了实现简洁与结果可复现。

\subsection{研究现状小结}
综合上述四条线索可以看出，现有研究在波形设计、波束赋形与轨迹优化三个层面均已取得丰富成果，但仍存在三处明确空白。

第一，\textbf{评价基准不统一}。加权和目标因量纲混杂与极值退化而难以横向比较，导致不同文献报告的“折中增益”无法直接对照。第二，\textbf{感知度量标度存在系统性误用}。CRB 对阵元数的三次方标度被广泛套用，而其成立前提（单阵元信噪比定义）在系统级建模中往往不满足。第三，\textbf{可行性与不可行点的报告不充分}。大量工作只绘制可行区内的折中曲线，对可行域边界、非单调点与未实现模块缺乏明确标注，使得读者难以判断方法的真实适用边界。

本文的三项创新正是针对这三处空白：以等速率约束替代加权和、给出并验证正确的 CRB 标度、以及在报告中如实标注全部边界与局限。

\section{本文主要研究内容与创新点}
\label{sec:contribution}

针对低空智联网复杂空域环境下的通感资源竞争问题，本文开展面向低空无人机通感一体化系统的理论建模、等速率感知性能优化与轨迹协同研究，主要工作与创新点如下：

\begin{enumerate}
  \item \textbf{构建低空无人机空地多载波通感一体化系统模型并推导角度估计 CRB 闭式解}：统筹考虑无人机三维几何拓扑、莱斯衰落信道与雷达回波散射特性，推导出均匀线阵角度估计克拉美罗下界（CRB）的闭式表达式
  \begin{equation}
    \mathrm{CRB} = \frac{6\sigma^{2}}{\pi^{2}\cos^{2}\theta\,\bigl(N_{\mathrm{R}}^{2}-1\bigr)\lvert\alpha\rvert^{2}N_{\mathrm{sym}}\,G(\theta)},
    \label{eq:intro_crb}
  \end{equation}
  明确指出其随 $N_{\mathrm{R}}^{-2}$ 而非 $N_{\mathrm{R}}^{-3}$ 标度，并解释文献中三次方表述源于单阵元信噪比定义这一前提差异。第四章的蒙特卡洛实验以五个阵元数配置验证闭式解，实测均方根误差与闭式解的比值落在 $[0.951,\,1.050]$ 区间内，最大偏差 5.0\%。

  \item \textbf{提出等通信速率约束下最小化感知 CRB 的联合波束赋形与子载波分配方法}：以匹配滤波感知波束、迫零/最小均方误差通信波束与一维闭式功率分配构成可行解，避免对不可实现的半正定松弛的依赖；证明了“通信功率取可行下界、剩余功率全部用于感知”这一结构的最优性；并利用交换论证证明按“单位功率通信效率”排序选取感知子载波的最优性，将规模为 $\binom{M}{n_s}$ 的组合搜索降为 $O(M\log M)$ 的一次排序。

  \item \textbf{建立“等 CRB 速率增益”的公平对比框架}：以相同感知 CRB 为基准比较各方案可支撑的通信速率，而非比较加权和目标值。该框架揭示联合优化相对时隙划分在等速率下的感知优势为 1.28$\times$–1.94$\times$，且该优势随目标速率提高单调扩大；同时如实报告 $k=48$ 处因感知子载波与通信子载波结构性重合而出现的 0.00\% 增益点，以及 $k=63$ 处回升至 4.95\% 的非单调现象。

  \item \textbf{提出速度约束下的无人机轨迹优化并给出可复现的诚实评测规范}：以序列二次规划求解逐时隙轨迹，相对绕目标匀速盘旋基线将稳态时隙的角度估计 CRB 降低约 52 倍；全部仿真脚本、参数与结果数据随论文一并提供，对不可行点（$\gamma\ge 8$\,dB）、非单调点、轨迹初始条件不对称以及未实现模块（半正定松弛、扩展卡尔曼滤波、学习类基线）逐一如实标注。
\end{enumerate}

需要说明的是，本文方法并非 $\mathcal{P}_{1}$ 的全局最优解。第三章将明确论证问题的非凸性来自速率约束的对数行列式结构与 CRB 对波束的双线性依赖两处，并说明本文放弃“先松弛再恢复”路线、转而构造可行解的理由：解一定可行、一定可复现、每一步物理含义可解释。这一取舍本身是本文方法论立场的一部分——在感知与通信的权衡研究中，一个可验证的可行构造比一个无法复现的“最优解”更有工程价值。

\section{论文组织架构}
\label{sec:organization}

本文共分为五章，具体结构安排如下：

第一章为绪论，阐述低空智联网与无人机通感一体化的研究背景、五项关键理论矛盾、国内外研究现状的四条线索及本文主要贡献。

第二章为国内外研究现状与理论基础，详细梳理通感一体化三大演进体制、空地信道特性与数学基准模型，为后续章节提供符号与工具准备。

第三章为核心理论与方法，依次给出系统模型与 CRB 闭式推导、等速率 CRB 最小化问题建模与可行域分析、波束赋形与功率分配（含功率结构最优性证明）、子载波选择最优性证明（含交换论证与近似条件讨论）、轨迹优化模型（含求解器选择与数值稳定性）及复杂度分析。

第四章为仿真评估，给出仿真参数、CRB 闭式解的蒙特卡洛验证、波束方向图、等速率感知性能对比、子载波分配权衡、PAPR 抑制效果、算法收敛性与轨迹优化结果，并设专节讨论轨迹对比的公平性与本文方法的六条适用边界。

第五章为总结与展望，归纳全文研究成果，指出本文方法在宽带斜视、多目标多径、非理想 CSI、有限字符与射频损伤、硬件验证与学习类基线六个方向上的可拓展空间。
